\documentclass[11pt]{article}
\usepackage[margin=1in]{geometry}
\usepackage{amsmath,amssymb,amsthm,mathtools}
\usepackage{graphicx}
\usepackage{microtype}
\usepackage[hidelinks]{hyperref}
\usepackage{enumitem}

\newtheorem{theorem}{Theorem}
\newcommand{\C}{\mathbb C}
\newcommand{\cH}{\mathcal H}

\title{A Square-Zero Counterexample to Conjecture 2.1\\
in \emph{A Conditional Proof of the ISP for Quasinilpotent Operators}}
\author{Automated mathematics run \texttt{fa\_banach\_001}, lane 16}
\date{August 11, 2026}

\begin{document}
\maketitle

\begin{center}
\fbox{\parbox{0.9\textwidth}{\textbf{Status.}
Candidate counterexample, likely valid, pending expert review. Conjecture 2.1
of Norman's arXiv:2008.03253v2 is false as printed. The construction is a
rank-one square-zero operator supported on a two-dimensional subspace of the
required infinite-dimensional Hilbert space.}}
\end{center}

\section{Source conjecture}

Norman~\cite{Norman} works on a complex separable infinite-dimensional
Hilbert space. His Conjecture 2.1 begins with a quasinilpotent operator $T$
and asserts the existence of $a,b,t>0$. Write $\widetilde T=mT$, where
$m>0$ is chosen so that $\|\widetilde T\|=a$. For every rank-one
quasinilpotent $F$ of norm $b$ such that
$\widetilde T^*+\alpha F$ is nonquasinilpotent for at least two distinct
complex values of $\alpha$, the conjecture requires a set
$S\subset\C\setminus\{0\}$, $|S|\ge2$, such that:
\begin{enumerate}[label=(\roman*)]
\item $\widetilde T^*+\alpha F$ is nonquasinilpotent for all $\alpha\in S$;
\item for every $\alpha\in S$,
\[
 \sigma_t(\widetilde T+\alpha F)
 \subseteq
 \sigma(\widetilde T+\alpha F)
   +B\!\left(0,\Phi_{a,b,t,S}(\alpha)\right),
\]
and the set on the right is disconnected in $\C$.
\end{enumerate}
The counterexample below uses the fact that the hypothesis is imposed on the
adjoint perturbation while the disconnectedness conclusion is imposed on the
non-adjoint perturbation.

\begin{figure}[p]
  \centering
  \includegraphics[width=0.86\textwidth]{figures/open_problem_crop.png}
  \caption{Conjecture 2.1 in Norman~\cite[p.~5]{Norman}. The full readable
  page width is retained so that the adjoint in the hypothesis and its absence
  in conclusion (ii) can be checked directly.}
  \label{fig:question}
\end{figure}

\section{Proof intuition}

The adjoint asymmetry can already be maximized by the smallest nilpotent
shift. Let $N$ send one basis vector to another and vanish everywhere else.
Perturbing $N$ by a scalar multiple of itself never changes square-zero
nilpotence, so the non-adjoint spectrum in conclusion (ii) stays equal to
$\{0\}$. In contrast, adding the same $N$ to $N^*$ closes a two-cycle. The
resulting $2\times2$ block has two nonzero eigenvalues for every nonzero
perturbation parameter. Thus the conjecture's hypothesis holds as strongly as
possible, while the required thickened spectrum is always a disk and can
never be disconnected.

\section{Counterexample}

\begin{theorem}\label{thm:counterexample}
Conjecture 2.1 of arXiv:2008.03253v2 is false as stated.
\end{theorem}

\begin{proof}
Let $\cH$ be a complex separable infinite-dimensional Hilbert space and fix
orthonormal vectors $e_1,e_2$. Define $N\in B(\cH)$ by
\[
 Ne_2=e_1,\qquad Ne_1=0,\qquad
 N|_{\{e_1,e_2\}^{\perp}}=0.
\]
Then $N$ has rank one, $N^2=0$, and $\|N\|=1$. In particular, $T:=N$ is
quasinilpotent.

We show that no triple $a,b,t>0$ can satisfy the conjecture for this $T$.
Fix arbitrary $a,b,t>0$. Since $\|T\|=1$, the normalization specified in
the conjecture is
\[
 \widetilde T=aN.
\]
Choose
\[
 F=bN.
\]
Then $F$ has rank one, $F^2=0$, and $\|F\|=b$, so it is one of the
perturbations quantified over by the conjecture.

Let $E=\operatorname{span}\{e_1,e_2\}$. Relative to the ordered basis
$(e_1,e_2)$, the adjoint-side perturbation restricts to
\[
 (\widetilde T^*+\alpha F)|_E
 =aN^*+\alpha bN
 =\begin{pmatrix}0&\alpha b\\ a&0\end{pmatrix}.
\]
Its square is
\[
 \begin{pmatrix}0&\alpha b\\ a&0\end{pmatrix}^{\!2}
 =\alpha ab I_E.
\]
Consequently, for every $\alpha\ne0$ the restriction has the two nonzero
eigenvalues satisfying $\lambda^2=\alpha ab$. Hence
$\widetilde T^*+\alpha F$ is not quasinilpotent for every nonzero $\alpha$.
In particular, the premise of Conjecture 2.1 holds for more than two distinct
values of $\alpha$.

On the non-adjoint side, however,
\[
 \widetilde T+\alpha F=(a+\alpha b)N,
 \qquad
 (\widetilde T+\alpha F)^2=0.
\]
Therefore
\[
 \sigma(\widetilde T+\alpha F)=\{0\}
 \quad\text{for every }\alpha\in\C.
\]
For any admissible ball radius $r$, its spectral thickening is
\[
 \sigma(\widetilde T+\alpha F)+B(0,r)=B(0,r),
\]
which is connected. If the open-ball convention $B(0,0)=\varnothing$ is
used at radius zero, the required inclusion of the nonempty pseudospectrum
already fails; under the singleton convention the thickening is still
connected. Thus condition (ii) of the conjecture fails for every nonzero
$\alpha$, regardless of the choice of the function $\Phi$.

The hypothesis and condition (i) hold for all nonzero $\alpha$, but no set
$S$ with $|S|\ge2$ can satisfy condition (ii). Since $a,b,t>0$ were
arbitrary, no existential triple works for $T=N$. This disproves the
conjecture.
\end{proof}

\section{Consequences and limitations}

Theorem~\ref{thm:counterexample} answers the source conjecture negatively and
blocks the paper's conditional deductions insofar as they rely on Conjecture
2.1 exactly as printed. It does not solve the Hilbert-space invariant
subspace problem. It also does not address a modified statement in which the
hypothesis and the pseudospectral conclusion use the same adjoint
orientation; the present construction is tailored to the printed asymmetry.

The value of $t$ and the detailed definition of the radius function $\Phi$
play no role: once the non-adjoint spectrum is $\{0\}$, every metric
thickening of it is connected. The example is finite rank but acts on the
infinite-dimensional space required by the source, by extending the
$2\times2$ block by zero.

\section{Verification and novelty}

The proof audit in \path{VERIFICATION.md} checks the source quantifiers,
rank and norm conditions, both perturbation orientations, and the final
connectedness obstruction. The script \path{code/check_matrices.py} checks
36 numerical instances of the two displayed matrix identities; it is only a
sanity check, not part of the proof.

The local run indexes were searched for arXiv:2008.03253, the exact title,
rank-one quasinilpotent perturbations, and invariant-subspace terms. Bounded
web searches on August 11, 2026 used the exact source title together with
``counterexample,'' the arXiv id together with ``Conjecture 2.1,'' and the
author together with the quasinilpotent invariant-subspace conjecture. They
found the current source version and adjacent invariant-subspace papers, but
no later answer or this counterexample. The official arXiv record lists v2,
revised August 25, 2021, as current. Novelty confidence is positive but
conservative pending expert review.

Human review is strongly recommended. The decisive audit is the literal
adjoint orientation in Figure~\ref{fig:question}; after that, the proof is
elementary two-dimensional operator algebra.

\begin{thebibliography}{9}
\bibitem{Norman}
M. Norman,
\emph{A conditional proof of the ISP for quasinilpotent operators},
arXiv:2008.03253v2 [math.FA], 2021.
\end{thebibliography}

\end{document}
